% Chapter 1 --- Multimodal prediction of Barrett's oesophagus progression
% Draft generated from the scientific-hardening analysis (branch chapter1-scientific-hardening).
% Figures embedded by artifact reference; compile within the thesis template.
% Citations use \cite keys resolved in chapter1_refs.bib --- VERIFY every DOI before submission.

% FIGURES: place the 6 PNG artifacts in a ./figures/ subdirectory next to this .tex.
% Filenames used below (download from the artifact store):
%   consort_diagram.png, fig_pr_fusion_vs_single.png, fig_robustness_forest.png,
%   fig_calibration_curves.png, fig_decision_curves.png,
%   fig_label_noise_attention.png, fig_encoder_sensitivity.png

\chapter{Multimodal prediction of progression in Barrett's oesophagus}
\label{ch:multimodal-barretts}

\section{Introduction}
\label{sec:intro}

Barrett's oesophagus is the principal precursor to oesophageal adenocarcinoma
(OAC), a cancer whose prognosis remains poor because most cases present at an
advanced stage. Endoscopic surveillance with biopsy aims to detect progression
early --- from non-dysplastic Barrett's (NDBE) through low-grade dysplasia (LGD)
to high-grade dysplasia (HGD) and invasive cancer --- but surveillance is
resource-intensive and histopathological grading of dysplasia suffers from
substantial inter-observer variability \cite{killcoyne2020}. Two data modalities
are routinely generated during surveillance yet rarely combined
computationally: shallow whole-genome sequencing (sWGS) of biopsy material,
which yields genome-wide copy-number variation (CNV) profiles, and
haematoxylin-and-eosin (H\&E) histopathology of the same biopsies.

This chapter asks a single, deliberately narrow question: \emph{does combining
CNV and H\&E histopathology improve prediction of near-term progression over
either modality alone, and is any such improvement robust, clinically
meaningful, and stronger than routine clinical information?} The emphasis is not
on novel architectures but on establishing whether a multimodal signal is real
and defensible at the sample size available. The chapter therefore reports a
frozen primary model together with an extensive programme of pre-specified
robustness, calibration, error, and sensitivity analyses.

\paragraph{Contributions.}
\begin{enumerate}
  \item A leakage-safe, patient-level evaluation of CNV-only, image-only, and
    fusion models for next-biopsy progression, with an independently
    reconstructed and verified endpoint.
  \item Evidence that late-mean CNV$+$histology fusion is the best-performing
    single model and improves discrimination and calibration over CNV alone,
    with the improvement stated honestly against each metric's confidence
    interval.
  \item A robustness programme showing the multimodal benefit strengthens on
    higher-provenance data subsets, is stable to aggregation and temporal
    choices, and is not driven by any single patient --- while documenting
    genuine fold-level heterogeneity.
  \item A systematic error analysis and interpretability investigation
    identifying \emph{tissue sampling} (a patient-level label applied to a
    single biopsy slide), rather than model artefact, as the dominant source of
    image-model error --- and a concrete structural argument for multimodality.
  \item Pre-specified encoder-sensitivity analyses (Virchow2, GigaPath, and
    multi-encoder stacking) establishing that the finding is not an artefact of
    one foundation model.
\end{enumerate}

\section{Background and related work}
\label{sec:background}

\subsection{Oesophageal adenocarcinoma and Barrett's surveillance}
Oesophageal adenocarcinoma (OAC) is an aggressive malignancy that is usually
diagnosed late and carries poor survival \cite{lagergren2017oac}. Its principal
precursor, Barrett's oesophagus, is subject to endoscopic surveillance with the
goal of detecting dysplastic progression early enough to intervene; current
gastroenterology guidelines codify biopsy protocols and surveillance intervals
\cite{shaheen2022acg,weusten2023esge}. Two limitations motivate computational
risk prediction. First, histopathological grading of dysplasia --- the current
basis for risk stratification --- suffers from substantial inter-observer
variability, and most surveillance biopsies are non-dysplastic even in patients
who later progress. Second, non-endoscopic triage tools such as the
Cytosponge-TFF3 test have broadened case-finding \cite{fitzgerald2020cytosponge}
but do not themselves stratify progression risk. Molecular and tissue-based
predictors have been developed to fill this gap: abnormal TP53 status predicts
progression independently of a dysplasia diagnosis \cite{redston2022tp53};
tissue-systems-pathology tests risk-stratify low-grade-dysplasia patients and
can outperform pathology review \cite{frei2020tissuecypher,khoshiwal2023tissuecypher};
and, most directly relevant to this chapter, genome-wide copy-number variation
from shallow whole-genome sequencing predicts progression years before
transformation \cite{killcoyne2020}. The present work builds on this
copy-number signal and asks whether histopathology, read computationally, adds
to it.

\subsection{Deep learning on histopathology}
Computational pathology has moved from task-specific convolutional networks
towards weakly supervised and foundation-model approaches. Early clinical-grade
systems demonstrated that whole-slide diagnosis could be learned from slide-level
labels alone \cite{campanella2019}, and that morphology-based models could
predict outcomes and molecular status directly from H\&E --- survival in
colorectal cancer \cite{kather2019survival}, microsatellite instability in
gastrointestinal cancer \cite{kather2019msi}, driver mutations in lung cancer
\cite{coudray2018lung}, and Gleason grade in prostate cancer
\cite{bulten2020gleason}. Because whole-slide images are gigapixel-scale and
annotations are sparse, the dominant paradigm is multiple-instance learning
(MIL): a slide is a bag of tile embeddings and a permutation-invariant
aggregator produces the slide-level prediction. Attention-based MIL
\cite{ilse2018abmil} and its clustering-constrained extension CLAM
\cite{lu2021clam} are the standard aggregators, and the attention weights offer
a route to interpretability. This chapter uses attention-based MIL as its
histology head.

\subsection{Foundation models as tile encoders}
The tile-embedding step is increasingly delegated to pathology foundation
models --- self-supervised encoders pretrained on very large slide corpora.
Transformer-based contrastive pretraining (CTransPath) \cite{wang2022ctranspath}
was an early demonstration; more recent general-purpose encoders such as UNI
\cite{chen2024uni}, the whole-slide model Prov-GigaPath \cite{xu2024gigapath},
and Virchow \cite{vorontsov2024virchow} are trained on hundreds of thousands of
slides and transfer broadly across tasks. This chapter uses UNI2 as its default
encoder and reports GigaPath and Virchow2 as pre-specified sensitivity analyses
(Section~\ref{sec:encoder}); a recurring question in the field, examined
directly here, is whether a larger or newer encoder reliably improves downstream
performance at small clinical sample sizes.

\subsection{Multimodal integration of histology and genomics}
Combining histopathology with molecular data is an active area, on the
hypothesis that morphology and genomics carry complementary prognostic
information. Pathomic Fusion introduced a framework for fusing histology and
genomic features for cancer prognosis \cite{chen2020pathomic}, and PORPOISE
extended integrative histology--genomic modelling across many cancer types
\cite{chen2022porpoise}; recent reviews survey the design space of multimodal
data integration in oncology \cite{lipkova2022multimodal}. Two themes from this
literature frame the present chapter. First, fusion strategy matters: late
(decision-level) fusion, early (feature-level) fusion, and intermediate
(learned joint) fusion trade off flexibility against the risk of over-fitting,
which is acute at the sample sizes typical of a single-cohort clinical study.
Second, much of this work is developed on large multi-cohort resources such as
TCGA; the added value of multimodality on a smaller, single-institution
surveillance cohort --- and whether it survives honest, leakage-controlled
evaluation --- is far less established. That is the gap this chapter addresses
for Barrett's oesophagus, where, to our knowledge, CNV and computational
histopathology have not previously been combined for progression prediction.


\section{Data and cohort}
\label{sec:cohort}

\subsection{Cohort construction}
The analysis cohort comprises \textbf{150 patients} contributing \textbf{707
pre-event biopsy sample-rows} (359 distinct biopsies; multiple sequenced
sample-rows may map to one physical biopsy). Fifty patients (33.3\%) are
progressors. Each row corresponds to one surveillance biopsy at which both a CNV
profile and an H\&E whole-slide image are available, taken \emph{before} the
event that defines the label. Patients are partitioned into five
patient-disjoint cross-validation folds, so that no patient contributes rows to
both training and test in any fold (Figure~\ref{fig:consort}).

\begin{figure}[htbp]
  \centering
  \includegraphics[width=0.85\textwidth]{figures/consort_diagram.png}
  \caption[Cohort consort diagram]{\textbf{Cohort construction.} 150 patients
  $\rightarrow$ 707 pre-event biopsy sample-rows (359 distinct biopsies), split
  by outcome (50 progressors / 100 non-progressors) and evaluated under 5-fold
  patient-disjoint cross-validation. Counts are patient-level unless stated as
  sample-rows.}
  \label{fig:consort}
\end{figure}

\subsection{Primary endpoint contract}
\label{sec:endpoint}
The primary endpoint is \texttt{NextBiopsyProgression\_LGD2plus}, a Boolean
defined per pre-event biopsy sample-row. A row is \textbf{positive} if the
patient's \emph{next} biopsy is HGD, intramucosal carcinoma, or OAC (grade
$\ge 3$), \emph{or} is LGD (grade 2) in a patient who already has a prior LGD in
their history (i.e.\ a second consecutive LGD); a single isolated LGD is not an
event. The event biopsy must occur at a strictly later date than the index
biopsy.

To guard against silent label drift, the endpoint was independently
reconstructed from source columns (next-biopsy grade and LGD-streak history)
using the locked rule. The reconstruction reproduced the stored endpoint on
\textbf{all 707 rows with zero discrepancies} (107 positive / 600 negative
sample-rows; 50 / 150 positive patients), and all 707 next-biopsy dates were
confirmed strictly later than the index date (no same-day events). The cohort is
termed \emph{pre-high-grade} rather than dysplasia-free, because prevalent LGD
and indeterminate-grade biopsies are present at index.

\section{Methods}
\label{sec:methods}

\subsection{Copy-number model}
Genome-wide relative copy-number features (arm-level and fixed-window
log$_2$ ratios) were derived from sWGS. The CNV model is a random forest
(500 trees, maximum depth 20) applied to a PCA-reduced feature representation,
with a configuration fixed in advance (\texttt{cnv\_rf\_conservative}). CNV
features are backbone-independent and identical across all image-encoder
experiments.

\subsection{Histopathology model}
Each whole-slide image is tiled and encoded with a frozen pathology foundation
model. The default encoder is UNI2 \cite{chen2024uni}; GigaPath
\cite{xu2024gigapath} and Virchow2 \cite{zimmermann2024virchow2} are used in
sensitivity analyses (Section~\ref{sec:encoder}). Per-slide tile embeddings form
a bag that is aggregated by an attention-based multiple-instance-learning
(ABMIL) head \cite{ilse2018abmil} (hidden dimension 256, attention dimension
128, dropout 0.1), trained to predict the slide-level label. Because the
encoders are frozen and applied per slide independently of any fold or label,
they introduce no train/test leakage.

\subsection{Fusion}
The headline fusion is \textbf{late-mean}: the per-sample probability average of
the CNV-only and image-only models, evaluated per held-out fold. A trained joint
CNV$+$image model (intermediate fusion) and clinical-stacking variants are also
reported. Late-mean was chosen as the primary fusion because it has no learned
fusion parameters and therefore cannot overfit the fusion step at this sample
size.

\subsection{Clinical baselines}
Seven prediction-time-safe clinical features were used (current grade, current
grade normalised, LGD streak so far, maximum pathology so far, biopsy index,
days since previous biopsy, first-biopsy flag). Clinical baselines
(grade-only, history-only, grade$+$history, and a small tree) and clinical$+$modality
stacks were fitted with leakage-safe nested logistic regression / cross-fitted
stacking on the frozen folds.

\subsection{Evaluation protocol}
All models are evaluated at the \textbf{patient level}: sample-row probabilities
are aggregated to a patient score by the maximum (patient-max), which exactly
reproduces the locked headline numbers. Metrics are area under the
precision--recall curve (AUPRC), area under the ROC curve (ROC~AUC), and the
Brier score, reported together. Confidence intervals and paired model contrasts
use the paired bootstrap (2000 resamples) over patients. Thresholds, calibration
(Platt scaling \cite{platt1999}), and inner model selection use training folds
only; the held-out test fold is never used for selection.

\paragraph{Reporting rules.} Following a pre-registered set of statistical
rules, every result reports patient and positive-patient counts; AUPRC is
retained as a central metric and is never described as conclusive when its
confidence interval crosses zero; internal cross-validation is never described
as external validation; and negative findings are preserved.

\section{Primary results}
\label{sec:primary}

Table~\ref{tab:primary} reports all models at the patient level ($n=150$,
50 positive). Late-mean CNV$+$image is the best single model on AUPRC (0.630)
and Brier (0.184); the full clinical$+$CNV$+$image stack achieves the best
ROC~AUC (0.801).

\begin{table}[htbp]
  \centering
  \caption[Primary model comparison]{\textbf{Primary model comparison}
  (patient-level, $n=150$, 50 positive patients). Best AUPRC/Brier: late-mean
  fusion. Best ROC~AUC: clinical$+$CNV$+$image.}
  \label{tab:primary}
  \begin{tabular}{lccc}
    \hline
    Model & AUPRC & ROC AUC & Brier \\
    \hline
    Prevalence (reference)        & 0.333 & 0.500 & 0.255 \\
    Clinical (grade $+$ history)  & 0.495 & 0.720 & 0.235 \\
    CNV only                      & 0.538 & 0.663 & 0.216 \\
    Image only                    & 0.557 & 0.731 & 0.245 \\
    CNV $+$ image (late-mean)     & \textbf{0.630} & 0.774 & \textbf{0.184} \\
    CNV $+$ image (joint MLP)     & 0.562 & 0.746 & 0.245 \\
    Clinical $+$ CNV              & 0.528 & 0.740 & 0.244 \\
    Clinical $+$ image            & 0.554 & 0.778 & 0.228 \\
    Clinical $+$ CNV $+$ image    & 0.576 & \textbf{0.801} & 0.222 \\
    \hline
  \end{tabular}
\end{table}

The pre-specified principal contrast is late-mean fusion versus CNV-only
(Figure~\ref{fig:prcurves}). The improvement is \textbf{conclusive on ROC~AUC
and Brier but not on AUPRC}: $\Delta$ROC~AUC $=+0.111$ (95\% CI $+0.001$ to
$+0.221$, excludes zero); $\Delta$Brier $=-0.032$ (95\% CI $-0.060$ to $-0.003$,
excludes zero); $\Delta$AUPRC $=+0.091$ (95\% CI $-0.031$ to $+0.219$,
\emph{crosses zero}). Per the reporting rules the AUPRC gain is reported as
directionally consistent but not conclusive on the full cohort. The marginal
value of histology \emph{given} clinical$+$CNV information is also not conclusive
($\Delta$AUPRC $=+0.048$, 95\% CI $-0.025$ to $+0.124$).

\begin{figure}[htbp]
  \centering
  \includegraphics[width=0.8\textwidth]{figures/fig_pr_fusion_vs_single.png}
  \caption[Precision--recall, fusion vs single modality]{\textbf{Headline
  precision--recall.} CNV-only, image-only, and late-mean fusion on the primary
  endpoint (patient-level, $n=150$/50 positive). Late-mean fusion dominates both
  single modalities; the dotted line marks prevalence.}
  \label{fig:prcurves}
\end{figure}

\section{Robustness}
\label{sec:robustness}

The stability of the principal contrast (late-mean $-$ CNV-only $\Delta$AUPRC)
was probed across matching provenance, endpoint definition, temporal exclusion,
and patient-aggregation choices (Table~\ref{tab:robustness},
Figure~\ref{fig:forest}). The benefit behaves as the opposite of a data
artefact: it \emph{strengthens} on the cleanest data. Restricting to
exactly-matched rows raises $\Delta$AUPRC to $+0.176$ (95\% CI $+0.007$ to
$+0.341$, excludes zero), and on strong grade-provenance subsets to $+0.161$
(95\% CI $+0.012$ to $+0.318$). The full-cohort estimate is diluted by a
high-prevalence fuzzy-matched stratum, not inflated by it. The benefit
attenuates at the stricter LGD3$+$ endpoint ($+0.049$, crosses zero) --- an
honest limitation --- and is stable to a 90-day minimum horizon and to every
aggregation scheme.

\begin{table}[htbp]
  \centering
  \caption[Robustness of the multimodal benefit]{\textbf{Robustness of the
  late-mean $-$ CNV-only $\Delta$AUPRC.} Prevalence is reported for every
  restricted cohort (reporting rule). ``Excl.\ 0'' indicates the 95\% CI
  excludes zero.}
  \label{tab:robustness}
  \begin{tabular}{llcccc}
    \hline
    Axis & Cohort & $n$ rows & Prev. & $\Delta$AUPRC & Excl.\ 0 \\
    \hline
    Reference           & full cohort            & 707 & 0.151 & $+0.091$ & no \\
    Exact matching      & exact-only             & 532 & 0.096 & $+0.176$ & \textbf{yes} \\
    Grade provenance    & scraped-confirmed      & 547 & 0.113 & $+0.161$ & \textbf{yes} \\
    Grade provenance    & exclude master-fallback& 552 & 0.114 & $+0.152$ & \textbf{yes} \\
    Same-day exclusion  & (0 same-day events)    & 707 & 0.151 & $+0.091$ & no \\
    Temporal horizon    & min-horizon 90\,d      & 676 & 0.154 & $+0.085$ & no \\
    Cleaner endpoint    & LGD3$+$                & 707 & 0.113 & $+0.049$ & no \\
    Aggregation         & patient-mean           & 150 & ---   & $+0.070$ & no \\
    Aggregation         & most-recent            & 150 & ---   & $+0.111$ & no \\
    Aggregation         & first-eligible (16 pos)& 150 & ---   & $+0.053$ & no \\
    \hline
  \end{tabular}
\end{table}

\begin{figure}[htbp]
  \centering
  \includegraphics[width=0.85\textwidth]{figures/fig_robustness_forest.png}
  \caption[Robustness forest plot]{\textbf{Robustness forest.} Late-mean $-$
  CNV-only $\Delta$AUPRC with 95\% bootstrap intervals across sensitivity
  cohorts. The benefit strengthens under exact-matching and strong-provenance
  subsets (intervals excluding zero) and attenuates at the stricter LGD3$+$
  endpoint.}
  \label{fig:forest}
\end{figure}

\paragraph{Fold and patient influence.} The principal internal-validity caveat
is fold-level heterogeneity: the fold-wise $\Delta$AUPRC ranges from $-0.112$
(fold 1) to $+0.420$ (fold 2), with folds 3--5 at $+0.057$, $+0.057$, and
$+0.154$. Leave-one-patient-out analysis shows no single patient drives the
result (maximum single-patient influence 0.023 on a base of 0.091; 82 positive
and 65 negative influencers), so the fragility is fold-level, not outlier-level.

\section{Calibration and clinical utility}
\label{sec:calibration}

The frozen fusion and CNV models are well calibrated (fusion calibration slope
1.78, expected calibration error [ECE] 0.075; CNV slope 1.23, ECE 0.095),
whereas the clinical stacks are over-confident (slopes 0.5--0.9, ECE
0.21--0.23) (Figure~\ref{fig:calibration}, left). At an operating point of
sensitivity $\ge 90\%$, the clinical$+$CNV$+$image model achieves specificity
0.66, positive predictive value 0.57, and negative predictive value 0.93
(patient prevalence 0.33). Decision-curve analysis \cite{vickers2006dca}
(Figure~\ref{fig:calibration}, right) shows the frozen fusion model has the
highest net benefit across the clinically relevant threshold range.

\begin{figure}[htbp]
  \centering
  \includegraphics[width=0.48\textwidth]{figures/fig_calibration_curves.png}
  \hfill
  \includegraphics[width=0.48\textwidth]{figures/fig_decision_curves.png}
  \caption[Calibration and decision-curve analysis]{\textbf{Calibration (left)
  and decision-curve analysis (right).} The frozen fusion and CNV models are
  well calibrated; the clinical stacks are over-confident. The fusion model
  yields the highest net benefit across clinically relevant thresholds.}
  \label{fig:calibration}
\end{figure}

\section{Error analysis and interpretability}
\label{sec:error}

\subsection{Systematic error taxonomy}
Using a per-fold leakage-safe threshold (90th percentile of other-fold
negatives), patients fall into: correct non-progressor (88), missed progressor
(22), correct progressor (18), fusion-harm missed-progressor (10), fusion-harm
false-positive (9), and persistent false-positive (3). The two modalities agree
on 70.7\% of patients, where fusion accuracy is 0.76; on the 44 disagreement
cases fusion accuracy falls to 0.568 while at least one single modality was
correct in every case. Mean-fusion therefore averages away a correct unimodal
signal when the modalities conflict --- a concrete, pre-registered motivation for
learned gating fusion in future work. Missed progressors are more temporally
distant than caught ones (median 729 vs 548 days), arguing against short-horizon
leakage.

\subsection{Image label-noise and tissue sampling}
\label{sec:labelnoise}
A central question is whether low image scores on progressor patients represent
model failure or a benign biopsy of a patient whose dangerous tissue lies
elsewhere. Three lines of evidence point to the latter
(Figure~\ref{fig:labelnoise}). First, the image score rises monotonically with
the imaged biopsy's own grade (mean probability 0.34 for NDBE, 0.50 for
indeterminate, 0.52 for LGD), so the encoder reads the slide it is given.
Second, of 70 image-missed positive rows, 51\% have a \emph{benign} current
biopsy, and 63\% of all positive rows have a current grade below the patient's
maximum --- the dangerous tissue is on a different biopsy, so a low score is
correct for the sampled slide. Third, attention maps on five such cases land on
genuine epithelium, not on scanner artefacts, folds, or background. Crucially,
for the 36 image-missed benign-current rows the CNV model scores above baseline
(mean 0.212 vs 0.172), so CNV recovers field-level risk the sampled histology
cannot show --- a structural argument for multimodality.

A focus-metric quality-control check (variance-of-Laplacian on tissue tiles)
confirmed one flagged case as genuinely out-of-focus (an $\sim$18-fold lower
focus score than sharp slides) but also revealed that the bare metric conflates
optical blur with sparse/fragmented tissue; a robust blur QC must normalise for
tissue content. Blur affects only a small minority of slides and does not change
the primary conclusions.

\begin{figure}[htbp]
  \centering
  \includegraphics[width=\textwidth]{figures/fig_label_noise_attention.png}
  \caption[Image label-noise / tissue sampling]{\textbf{Why the image model
  ``misses'' progressors.} Attention overlays on five benign current biopsies of
  progressor patients. Attention lands on real epithelium with no focal
  high-grade lesion; the low image score is correct for the sampled tissue,
  because the progression-defining lesion is elsewhere in the patient. This is a
  limitation of the patient-level label, not model artefact, and CNV recovers
  the signal these slides cannot show.}
  \label{fig:labelnoise}
\end{figure}

\section{Encoder and multi-encoder sensitivity}
\label{sec:encoder}

To test whether the multimodal finding depends on the choice of pathology
foundation model, the primary endpoint was re-run with alternative encoders
under an identical architecture and protocol (Figure~\ref{fig:encoder}).
\textbf{Virchow2} performed slightly worse than UNI2 on every family and metric
(image-only AUPRC 0.482 vs 0.557, ROC~AUC 0.666 vs 0.731; late-mean AUPRC 0.558
vs 0.630); the ROC deltas exclude zero and Virchow2 was worse calibrated. Per
the pre-specified rule it was not expanded to other endpoints.

\textbf{GigaPath} was the strongest single image encoder on AUPRC (0.609), and a
fold-pure logistic stacker over the three encoders up-weighted GigaPath and
UNI2 and down-weighted Virchow2 --- so the encoders do carry different signal.
However, \emph{combining} encoders did not beat the single-encoder headline: the
three-encoder learned stack overfits at $n=150$ (AUPRC 0.535, significantly
worse than the headline, $\Delta=-0.095$, 95\% CI $-0.172$ to $-0.009$);
substituting GigaPath into the late-mean fusion ties the headline (0.624 vs
0.630); and fusing CNV with the mean of UNI2 and GigaPath gives the best point
estimate observed (0.635) but within noise. \textbf{UNI2 remains the reported
backbone} (with GigaPath an acceptable tie); the multimodal finding is not an
artefact of one foundation model.

\begin{figure}[htbp]
  \centering
  \includegraphics[width=0.7\textwidth]{figures/fig_encoder_sensitivity.png}
  \caption[Encoder sensitivity]{\textbf{Encoder sensitivity.} UNI2 versus
  Virchow2 on the primary task; UNI2 is equal or better on every family and
  metric. Multi-encoder stacking (Section~\ref{sec:encoder}) did not improve on
  the single-encoder headline.}
  \label{fig:encoder}
\end{figure}

\paragraph{Benchmark feasibility.} A subset of 69 of the 150 patients overlaps
the Killcoyne discovery cohort \cite{killcoyne2020}, but these patients are a
strict subset of the present training cohort and cannot serve as external
validation; moreover the published benchmark is CNV-only, precluding a
like-for-like multimodal comparison. The CNV-only leave-one-patient-out result
here (ROC~AUC 0.78--0.80) sits somewhat below the published $\sim$0.87, plausibly
reflecting preprocessing and QC differences.

\section{Discussion}
\label{sec:discussion}

Across a seven-part hardening programme, combining CNV and H\&E histopathology
by late-mean fusion is the best-performing single model for near-term
progression, improving ROC~AUC (0.66 $\rightarrow$ 0.77) and calibration
(Brier 0.216 $\rightarrow$ 0.184) over CNV alone, with a directionally
consistent AUPRC gain (0.54 $\rightarrow$ 0.63) whose full-cohort confidence
interval crosses zero but reaches significance on the cleanest data subsets. The
benefit is robust to aggregation and temporal choices and is not driven by any
single patient. The error analysis reframes the image model's apparent failures
as a consequence of patient-level labelling of individually-sampled biopsies,
and shows that CNV recovers exactly the field-level risk that a benign-appearing
slide cannot --- a mechanistic, not merely empirical, case for multimodality.

The strongest defensible statement is therefore: \emph{in an internally
validated, patient-disjoint cross-validation of 150 patients (50 progressors),
late-mean fusion of CNV and H\&E histopathology improved next-biopsy LGD2$+$
progression prediction over CNV alone in discrimination (ROC~AUC $+0.11$, 95\%
CI $+0.001$ to $+0.221$) and calibration (Brier $-0.032$, 95\% CI $-0.060$ to
$-0.003$); the AUPRC gain was directionally consistent and reached significance
on exactly-matched and strong-provenance subsets; the result is internal
validation only.}

\paragraph{Claims not made.} These results do not support describing the AUPRC
improvement as conclusive on the full cohort, describing any result as
externally validated or generalising, treating the Killcoyne overlap as external
validation, promoting a robustness subset to the headline, claiming histology
adds conclusive value over clinical$+$CNV, claiming a newer/larger encoder
improves performance, presenting sample-rows as independent patients, or
implying a deployment-ready screening tool from $n=150$.

\section{Limitations and future work}
\label{sec:limitations}

The study is \textbf{internally validated only}; no external cohort has been
tested, and predictive claims should not be extended beyond internal
cross-validation until a digitised external cohort is available. The patient-max
label imposes a structural floor of ``unwinnable'' cases on any biopsy-level
image model (Section~\ref{sec:labelnoise}); biopsy-level image labels and a
pathologist read of missed cases would quantify and reduce this label noise. The
fold-level heterogeneity (Section~\ref{sec:robustness}) indicates the point
estimate is sensitive to cohort composition at this sample size.

The ranked next investments are: (1) prepare the chapter and manuscript on these
hardened internal results; (2) obtain a digitised external cohort before making
further predictive claims; (3) add pathologist-supported interpretation of
disagreement cases; and (4) --- only if specifically required --- a clearly
labelled within-cohort Killcoyne benchmark. Learned gating fusion, motivated by
the modality-disagreement finding, is a natural direction for a subsequent
chapter but is deferred here to avoid overfitting the fusion step at the present
sample size.
